\chapter{Lista de recetas}
\label{chap:recetas}
Esta es la función de recetas de la aplicación.
Muestra una lista de recetas y permite añadir una nueva, editar una existente y eliminarla.

La lista y el almacenamiento de las recetas parten de lo visto en clase.
Se sigue la misma paleta de colores que en el resto de la aplicación.

\begin{figure}[H]
  \centering
  \includegraphics[width=0.4\textwidth]{recetas}
\end{figure}

\section{Lista de recetas}
La pantalla de la lista usa un \texttt{RecyclerView} con un adaptador que dibuja una tarjeta por cada receta.
Al pulsar una tarjeta se abre el detalle de esa receta, y el botón de nueva receta abre el formulario vacío.
Cuando se vuelve a la lista, se llama a \texttt{notifyDataSetChanged()} desde \texttt{onResume} para que se vean los cambios.

\section{Edición de una receta}
Cada receta tiene un botón Editar en su pantalla de detalle.
Reutiliza el formulario de nueva receta en lugar de crear una pantalla aparte.

El detalle envía la posición de la receta en el \texttt{Intent} y cierra su pantalla:

\begin{lstlisting}
findViewById<Button>(R.id.botonEditar).setOnClickListener {
    val intent = Intent(this, NuevaRecetaActivity::class.java)
    intent.putExtra(NUMERO_RECETA, idReceta)
    startActivity(intent)
    finish()
}
\end{lstlisting}

\section{Crear o editar con el mismo formulario}
El formulario lee la posición recibida.
Si no recibe ninguna, el valor por defecto es $-1$ y se crea una receta nueva.
Si recibe una posición de $0$ o más, se editan los datos de esa receta.

\begin{lstlisting}
idEdicion = intent.getIntExtra(RecetaActivity.NUMERO_RECETA, -1)

if (idEdicion >= 0) {
    rellenarFormulario(RepositorioSingleton.instance.listaRecetas.buscar(idEdicion))
} else {
    anadirFilaIngrediente()
    anadirFilaPaso()
}
\end{lstlisting}

La función \texttt{rellenarFormulario} cambia el título a Editar receta y el botón a Guardar cambios.
También rellena los campos, selecciona la categoría y crea una fila por cada ingrediente y por cada paso de la receta.
Un ingrediente sin cantidad se muestra como al gusto.

\section{Guardar los cambios}
Al guardar, el formulario construye una receta con los datos escritos.
Si está en modo edición llama a \texttt{actualizar(idEdicion, receta)}, que sustituye la receta de esa posición.
Si no, llama a \texttt{nueva(receta)}, que la añade al final.
Después cierra la pantalla y la lista se actualiza al volver.